Thus, to complete the proof, it suffices to note that:
@@M0@@

To show that @@M1@@ is zero, we take @@M2@@, and @@M3@@ so that there is a complex @@M4@@ such that @@M5@@, and @@M6@@; since @@M7@@ has dimension @@M8@@, this implies @@M9@@, hence @@M10@@ is acyclic, and @@M11@@.

\sgaStableTarget{sga6.e06.lem.6-7}\textbf{Lemma 6.7.} \uline{Let} @@M12@@ \uline{be an affine noetherian scheme}, @@M13@@ \uline{finitely many points of} @@M14@@, and @@M15@@ \uline{an invertible} @@M16@@\uline{-Module}. \uline{Then} \uline{there is a section} @@M17@@ \uline{such that, for every} @@M18@@, @@M19@@ \uline{is a basis of} @@M20@@.

Put @@M21@@. For every @@M22@@, let @@M23@@ be a closed point of @@M24@@ such that @@M25@@. The @@M26@@ correspond to maximal ideals @@M27@@ of @@M28@@, which we may suppose all distinct, changing the index set so that each of the @@M29@@ occurs once and only once in the family @@M30@@. Since @@M31@@ is invertible, for every @@M32@@ there is an isomorphism @@M33@@ from which we deduce an isomorphism @@M34@@, then @@M35@@, or again @@M36@@. Lifting this last isomorphism, we obtain a homomorphism @@M37@@ such that the corresponding homomorphism @@M38@@ is, for every @@M39@@, an isomorphism modulo @@M40@@, hence an isomorphism by Nakayama. The image of @@M41@@ under this homomorphism therefore gives a section @@M42@@ of @@M43@@ that is a basis of @@M44@@ for every @@M45@@. Since the set of points where @@M46@@ is a basis of @@M47@@ is an open subset of @@M48@@, @@M49@@ is also a basis of @@M50@@ for @@M51@@.

\sgaStableTarget{sga6.e06.lem.6-8}\textbf{Lemma 6.8.} \uline{Let} @@M52@@ \uline{be a noetherian scheme having an ample sheaf} @@M53@@, @@M54@@ \uline{an invertible} @@M55@@\uline{-Module}, and @@M56@@ \uline{points of} @@M57@@. \uline{Then there is an}
